\documentclass{article}

% NeurIPS 2026 — anonymous submission
\usepackage[preprint]{neurips_2026}

\usepackage[utf8]{inputenc}
\usepackage[T1]{fontenc}
\usepackage{hyperref}
\usepackage{url}
\usepackage{booktabs}
\usepackage{amsfonts}
\usepackage{amsmath}
\usepackage{amssymb}
\usepackage{nicefrac}
\usepackage{microtype}
\usepackage{xcolor}
\usepackage{graphicx}
\usepackage{subcaption}
\usepackage{multirow}
\usepackage{array}
\usepackage{natbib}

% ── Placeholder macro ─────────────────────────────────────────────────────────
\newcommand{\placeholder}[1]{%
  \colorbox{yellow!40}{\textbf{[#1]}}%
}

\title{Compositional Reasoning Depth Predicts Clinical AI Failure:\\
Empirical Evidence Consistent with the Peng Impossibility Theorem\\
in Electronic Health Record Question Answering}

\author{Anonymous Authors}

\begin{document}

\maketitle

% ── Abstract ──────────────────────────────────────────────────────────────────
\begin{abstract}
Aggregate accuracy benchmarks conceal a systematic structure in clinical AI failure:
questions requiring more inferential steps produce disproportionately more errors.
Motivated by the Peng impossibility theorem for transformer attention composition,
we introduce a pre-registered \emph{hop-count taxonomy}---the number of distinct
reasoning steps required to answer a clinical question from an electronic health
record (EHR)---as a principled predictor of model failure.
We annotate 313 clinician-generated EHR question-answer pairs from MedAlign
across four hop levels and evaluate 301 questions in a within-model ablation
(\texttt{claude-sonnet-4-6}, zero-shot versus extended thinking) and a
cross-architecture replication (\texttt{gpt-4o}, zero-shot).
Both models exhibit monotone accuracy decline with hop count:
Claude Sonnet zero-shot declines from 30.6\% (hop=1) to 17.3\% (hop=4)
(Cochran--Armitage $z = -2.38$, $p = 0.009$; OR per hop $= 0.72$,
95\% CI [0.56, 0.92], $p = 0.008$); GPT-4o independently replicates this
pattern (37.8\% to 14.7\%; CA $z = -3.80$, $p < 0.0001$; OR $= 0.58$
[0.45, 0.75], $p < 0.001$).
A pre-specified context sufficiency audit confirmed that higher-hop questions
are not differentially disadvantaged by EHR truncation (answerability
93--95\% at hops 2--4 vs.\ 79\% at hop=1), establishing that the
degradation reflects compositional reasoning difficulty.
Extended thinking did not significantly flatten the accuracy-versus-depth curve
across three reasoning conditions (internal extended thinking, structured
four-step CoT, and 5$\times$ larger thinking budget), and thinking-token usage
scaled with hop count ($r = 0.31$, $p < 0.0001$), consistent with the
predicted $O(k)$ computational requirement.
These findings establish hop count as a theoretically grounded, cross-architecture
predictor of clinical AI failure with immediate implications for deployment
risk stratification.
\end{abstract}

% ── 1. Introduction ───────────────────────────────────────────────────────────
\section{Introduction}

Artificial intelligence systems based on transformer architectures are now
embedded in clinical workflows ranging from automated chart summarization to
real-time diagnostic support~\citep{singhal2023large,mcduff2025}.
Despite rapidly improving benchmark performance on standardized medical
examinations~\citep{jin2020disease,hendrycks2021measuring}, these systems
exhibit inconsistent and sometimes dangerous failures when applied to
real-world electronic health record (EHR) tasks~\citep{noharm2025}.
A persistent challenge is that aggregate accuracy metrics conceal the
structure of failure: researchers and clinicians typically observe \emph{that}
a model fails without understanding \emph{why} or \emph{when} it will fail.

We propose that compositional reasoning depth---the number of distinct
inferential hops required to produce a correct answer---is a principled
predictor of clinical AI failure, grounded in a formal theoretical result.
\citet{peng2024limitations} prove an impossibility theorem for transformer
self-attention: a single attention layer with $H$ heads, head dimension $d$,
and numerical precision $p$ cannot solve the $(n, k)$-function-composition
problem when $n \log n > H \cdot d \cdot p$. Informally, once the context
length $n$ exceeds a threshold determined by the model's architectural
parameters, the attention mechanism lacks sufficient representational capacity
to chain multiple inferential steps within a single forward pass.
Three theoretical escape routes exist: multi-step chain-of-thought reasoning
(which effectively adds layers), retrieval-augmented generation, and
non-standard architectural choices such as state-space models or mixture-of-expert routing~\citep{zhou2023brainformers,yehudai2025compositional}.
While the theorem formally characterizes single-layer attention, practical transformers
stack many such layers; we therefore treat it as a theoretical framework motivating
\emph{qualitative} predictions about compositional difficulty rather than applying it
quantitatively to the multi-layer setting. The critical empirical question is whether
accuracy degrades with hop count and whether test-time reasoning provides relief, predictions
that hold regardless of whether the single-layer impossibility is the proximate mechanism.

Clinical EHR question answering is a natural testing ground for this theorem.
A clinician asking ``Has this patient ever had a statin prescribed?'' requires
a single retrieval operation (hop=1). Asking ``Did the patient have
an elevated creatinine within 48 hours of contrast administration?'' requires
retrieving contrast administration timestamps, retrieving creatinine results,
and computing their temporal overlap (hop=3). Asking ``Summarize all
specialist consultations over the past year and identify any unresolved
recommendations'' requires synthesizing heterogeneous EHR entries across
multiple domains and time points (hop=4). The Peng theorem predicts that these
structurally distinct tasks should exhibit qualitatively different error rates,
and that test-time reasoning---which increases effective computation depth
without changing model weights---should selectively benefit higher-hop tasks.

This paper makes five contributions.
First, we introduce a pre-registered hop-count taxonomy for clinical EHR
question answering and annotate 313 questions from MedAlign across four hop levels.
Second, we provide the first cross-architecture empirical evidence that hop count
predicts clinical AI failure: a within-model ablation on \texttt{claude-sonnet-4-6}
(zero-shot vs.\ extended thinking) and an independent replication on \texttt{gpt-4o}
both confirm monotone accuracy decline with compositional depth.
Third, we conduct a context sufficiency audit demonstrating that the hop--accuracy
relationship reflects compositional reasoning difficulty rather than differential
EHR truncation, directly ruling out the leading alternative confound.
Fourth, we provide mechanistic evidence by showing that thinking-token usage
scales with hop count ($r = 0.31$, $p < 0.0001$), consistent with the
$O(k)$ computational requirement predicted by the Peng theoretical framework.
Fifth, we characterize the \emph{structure} of failure by hop level,
finding that omission and hallucination errors have distinct hop-count profiles
and that structured decomposition converts hallucination to omission---
a clinically favorable trade-off.
Sixth, we show through a question-type sensitivity analysis and a truncation
robustness ablation (32k-character context) that the hop--accuracy relationship
is not confounded by question type composition or EHR context availability,
and strengthens after controlling for reasoning type (adjusted OR = 0.57, $p = 0.005$).

% ── 2. Background and Related Work ───────────────────────────────────────────
\section{Background and Related Work}

\paragraph{The Peng impossibility theorem.}
\citet{peng2024limitations} establish that for any single-layer transformer
attention mechanism, the $(n, k)$-function-composition problem---which asks
whether the model can compose $k$ arbitrary functions over a context of length
$n$---is unsolvable when $n \log n > H \cdot d \cdot p$.
This is not an empirical observation but a formal impossibility result derived
from communication complexity theory. The theorem applies to the softmax
self-attention operation as standardized in modern transformers and does not
depend on training procedure, model scale, or the specific clinical domain.
Critically, the theorem identifies test-time chain-of-thought reasoning as a
theoretically sound mitigation: by decomposing a $k$-hop problem into $k$
sequential single-hop steps, each within the computability threshold,
CoT reasoning should flatten the accuracy degradation curve.
\citet{yehudai2025compositional} provide complementary analysis showing that
recurrent architectures and chain-of-thought can solve compositional tasks
that standard transformers cannot, independently corroborating the Peng
escape route. \citet{zhou2023brainformers} demonstrate architectural
alternatives (non-uniform block structures) that can route around the
attention bottleneck, though at engineering complexity cost.

\paragraph{Multi-hop reasoning benchmarks.}
The NLP community has extensively studied multi-hop reasoning in open-domain
QA settings. HotpotQA~\citep{yang2018hotpotqa} provides 113K questions
requiring evidence synthesis from two Wikipedia paragraphs;
MuSiQue~\citep{trivedi2022musique} constructs up to 4-hop questions by
composing single-hop questions with explicit bridge entities;
2WikiMultihopQA~\citep{ho2020constructing} adds cross-document coreference
as a reasoning challenge.
Across these benchmarks, model accuracy declines consistently with the number
of required reasoning hops, a pattern attributed to error propagation and
attention capacity limitations~\citep{yang2018hotpotqa,trivedi2022musique}.
However, none of these studies has (a) grounded the hop–accuracy relationship
in a formal theorem, (b) tested the theoretical CoT mitigation as a function
of hop depth rather than on average, or (c) applied hop-count analysis to
\emph{clinical} EHR data, where question structures, error stakes, and
ground-truth generation differ fundamentally from Wikipedia-based QA.

\paragraph{Clinical AI benchmarks and EHR reasoning.}
MedQA~\citep{jin2020disease}, MedBench, and MMLU Medical~\citep{hendrycks2021measuring}
evaluate LLMs on medical licensing examinations, which involve closed-book
multiple-choice questions requiring medical knowledge but not EHR navigation.
Med-PaLM~\citep{singhal2023large} achieved expert-level performance on such
benchmarks but was not evaluated on compositional EHR reasoning tasks.
emrQA~\citep{pampari2018emrqa} provides a large EHR QA benchmark derived from
i2b2 relations, but questions are primarily single-hop slot-filling and include
no compositional depth taxonomy.
Temporal reasoning in clinical NLP---including tlink prediction and timeline
reconstruction tasks~\citep{bethard2016semeval}---documents similar depth-dependent
difficulty patterns in clinical text, but does not frame these as a function of
a formal impossibility result or test reasoning-model mitigations.
FDARxBench~\citep{xiong2026fdarxbench} evaluates LLMs on document-grounded regulatory
biomedical QA and shows that retrieval and long-context grounding remain bottlenecks
even in best-in-class models---consistent with our finding that the hop--accuracy
relationship is not resolved by simply providing more context.
MedAlign~\citep{wornow2024medalign} provides clinician-generated, open-ended
questions paired with EHR XML records and clinician reference answers, making
it the most ecologically valid available benchmark for testing compositional
EHR reasoning.
No prior study has stratified MedAlign performance by compositional depth,
nor applied hop-count analysis to any clinical EHR benchmark.

\paragraph{Chain-of-thought and reasoning models.}
\citet{wei2022chain} demonstrated that chain-of-thought prompting elicits
multi-step reasoning in sufficiently large models, producing gains on
arithmetic and commonsense tasks.
\citet{kojima2022large} showed that a zero-shot CoT trigger (``let's think
step by step'') is sufficient to induce reasoning benefits without few-shot
exemplars.
\citet{nye2021show} proposed scratchpads as a formal mechanism for
intermediate computation, providing theoretical grounding for CoT.
Reasoning-specialized models---including OpenAI o1~\citep{openai2024o1}
and DeepSeek R1~\citep{deepseekr1_2025}---extend these ideas through
reinforcement learning over extended thinking traces, achieving large
improvements on mathematical and scientific benchmarks.
Information-theoretic analyses of CoT token requirements provide theoretical
grounding for P3: \citet{merrill2024expressive} prove that transformers with
polynomially many CoT tokens can solve problems inaccessible to direct inference,
and that the required token count scales with reasoning depth---a formal
analog of the $O(k)$ computation scaling our experiments test empirically.
Critically, no study has measured CoT benefit \emph{as a function of hop depth}
in clinical settings, leaving the predicted depth $\times$ compute interaction
untested in medicine.

\paragraph{Safety and variance in clinical AI.}
\citet{noharm2025} document systematic safety failures of clinical LLMs,
finding that models hallucinate clinical facts at rates that vary by
query type, but do not stratify by compositional depth.
Prior work on automation bias and de-skilling in clinical
AI~\citep{goddard2012automation,lyell2017automation} has not examined whether
these effects are moderated by task complexity operationalized as hop count.
The \emph{structure} of failure---hallucination versus omission versus
reasoning error---as a function of compositional depth has not been
previously characterized.

% ── 3. Problem Formulation ────────────────────────────────────────────────────
\section{Problem Formulation}

\paragraph{The Peng theorem and clinical QA.}
Let $\mathcal{E}$ denote the EHR context of length $n$ tokens, and let
$\mathcal{Q}$ denote a clinical question requiring composition of $k$ functions
$f_1, \ldots, f_k$ over $\mathcal{E}$, where each $f_i$ retrieves or transforms
a specific subset of the EHR.
The $(n, k)$-function-composition problem asks a model to produce a correct
answer by sequentially applying these functions.
By the Peng theorem~\citep{peng2024limitations}, a single attention layer with
parameters $(H, d, p)$ cannot solve this problem when $n \log n > H \cdot d \cdot p$.
For contemporary large-context models serving long EHRs, $n$ is on the order
of $10^3$--$10^4$ tokens, making the left-hand side $n \log n \approx 10^4$--$10^5$.
Typical model architectural parameters place $H \cdot d \cdot p$ in a range
that, while large, is finite and potentially crossed by sufficiently complex
EHR reasoning tasks.

The theorem generates three testable predictions:
(P1) \emph{Monotone degradation}: Accuracy should decrease monotonically with
$k$ under standard (non-CoT) inference.
(P2) \emph{CoT flattening}: Chain-of-thought reasoning, by decomposing
$k$-hop problems into $k$ sequential 1-hop steps, should flatten the
accuracy-versus-hop curve, with greater absolute benefit at higher $k$.
(P3) \emph{Computation scaling}: The amount of reasoning computation (proxied
by thinking token count in extended-thinking models) should scale with $k$,
consistent with the theoretical requirement for $O(k)$ sequential computations.

\paragraph{Scope of the theoretical mapping.}
The Peng theorem formally characterizes the computational capacity of a \emph{single}
attention layer. Modern transformers stack $L \gg 1$ such layers, which can in principle
decompose multi-hop reasoning through depth rather than through width or precision.
We therefore interpret P1--P3 as qualitative predictions derived from the theorem as a
motivating theoretical framework---not as strict corollaries of the formal impossibility
result for the multi-layer setting. The empirical claims of this paper stand on their own:
that compositional depth predicts accuracy, that thinking-token usage scales with depth,
and that test-time reasoning does not statistically significantly flatten the degradation
curve at current compute budgets. These findings are consistent with the Peng-motivated
predictions but do not require the single-layer proof as their sole justification.

\paragraph{Hop-count taxonomy.}
We define a pre-registered four-level hop taxonomy for clinical EHR questions
(Table~\ref{tab:hop_taxonomy}). Hop=1 questions require retrieval of a single
fact from the EHR (e.g., ``What is the patient's last creatinine?'').
Hop=2 questions require composing exactly two EHR facts or applying a single
comparison or threshold operation (e.g., ``Did the patient's blood pressure
exceed 160~mmHg at any point this admission?''---retrieving blood pressure
readings and comparing to a threshold).
Hop=3 questions require three-step inference chains involving cross-domain
retrieval, temporal reasoning, or contextual comparison
(e.g., ``Are any additional vaccines recommended for this patient?''---requiring
extraction of vaccination history, patient demographics, and application of
immunization guidelines).
Hop=4 questions require synthesis or generation across multiple EHR components
(e.g., ``Summarize all laboratory work from the past year''---requiring
identification, filtering, and narrative synthesis of heterogeneous lab entries).

\begin{table}[t]
\centering
\caption{Pre-registered hop-count taxonomy for clinical EHR question answering,
with operational definitions and representative examples from MedAlign.}
\label{tab:hop_taxonomy}
\small
\begin{tabular}{@{}lp{4.5cm}p{5.5cm}@{}}
\toprule
\textbf{Hop} & \textbf{Operational Definition} & \textbf{Representative Example} \\
\midrule
1 & Direct retrieval of a single named fact from the EHR.
  & ``What is the patient's current allergy list?'' \\
\addlinespace
2 & Two-fact composition, threshold application, or binary temporal comparison.
  & ``Did the patient have an elevated WBC on admission?'' \\
\addlinespace
3 & Three-step reasoning chain; requires cross-domain lookup, temporal arithmetic, or contextual clinical guideline application.
  & ``Is this patient eligible for supplemental breast cancer screening based on current guidelines?'' \\
\addlinespace
4 & Multi-component synthesis or generation requiring integration of $\geq$4 EHR sources or time spans.
  & ``Summarize all specialist consultations and flag unresolved recommendations.'' \\
\bottomrule
\end{tabular}
\end{table}

% ── 4. Methods ────────────────────────────────────────────────────────────────
\section{Methods}

\paragraph{Dataset: MedAlign.}
MedAlign~\citep{wornow2024medalign} consists of 983 clinician-generated
instructions paired with EHR XML files from the Stanford STARR-OMOP database
and clinician reference answers. We use the publicly released
subset of 402 unique question-answer pairs for which clinician-reviewed model
responses and source EHR files are available. After restricting to
questions with a uniquely identifiable EHR file and valid clinician reference
response, 313 questions were annotated and eligible; of these, 301 had
complete model evaluations (API errors or timeouts accounted for 12 omissions),
yielding the final analytic sample
(hop=1: $n=111$, hop=2: $n=46$, hop=3: $n=42$, hop=4: $n=104$;
602 total evaluations across both conditions).
EHR records were parsed from XML, chronologically structured, and truncated
to 8,000 characters (approximately 1,600--2,000 tokens at 4--5 characters per token,
well within the model's 200K-token context window) while retaining
the full character count as a completeness covariate.

\paragraph{Context sufficiency analysis.}
To test whether EHR truncation differentially disadvantaged higher-hop questions---
a potential confound whereby apparent compositional failure reflects information
unavailability rather than reasoning incapacity---we computed \emph{evidential coverage}
for each of the 301 questions. We extracted non-stopword content tokens ($\geq$4 characters)
from each clinician reference answer and computed the fraction that appeared verbatim in
the truncated 8,000-character EHR. We defined a question as ``any-token supported'' if
at least one reference content token appeared in the EHR, and ``well-supported'' if
$\geq$30\% of reference tokens appeared. A question where the answer evidence is absent
from the truncated EHR is implicitly unanswerable regardless of the model's reasoning
capacity; if higher-hop questions were systematically less well-supported, this would
constitute a truncation confound.

\paragraph{Hop annotation.}
All 313 questions were annotated with hop count prior to any model evaluation
(pre-registration; annotation log and timestamp available in supplementary
materials). Each question was presented to an LLM annotator together with a
structured rubric defining each hop level, with explicit examples and a
chain-of-thought elicitation prompt. To assess annotation reliability, a second independent annotation pass was
run on a stratified random sample of 50 questions (12--13 per hop level)
using the same model and prompt (within-model reproducibility), and a
cross-model validation was performed using \texttt{gpt-4o}.
Within-model reliability was substantial ($\kappa = 0.787$, 95\% CI
[0.603, 0.970]; 84\% exact agreement, 98\% agreement within $\pm$1 hop).
Cross-model agreement was moderate ($\kappa = 0.464$, 95\% CI [0.292,
0.637]; 60\% exact), reflecting divergent application of the rubric between
model families. All annotations were finalized and locked before any
model inference was run, precluding post-hoc adjustment.

\paragraph{Experimental design.}
We evaluate two conditions using \texttt{claude-sonnet-4-6}, a publicly
available large language model with extended-thinking capability:
\begin{itemize}
  \item \textbf{Condition A (Zero-Shot):} The model receives a patient EHR
    and question with the prompt ``Answer concisely'' and no explicit
    reasoning instructions. This is the dense baseline.
  \item \textbf{Condition B (Extended Thinking):} The same model receives the
    same prompt with extended thinking enabled
    (\texttt{thinking=\{"type":"enabled","budget\_tokens":3000\}}),
    allowing an internal chain-of-thought scratchpad before producing
    the visible response.
\end{itemize}
This within-model ablation design isolates the effect of test-time
reasoning from confounds such as model capacity, training data, or
architectural differences. Both conditions use identical weights, identical
prompts, and identical context. The sole manipulation is whether reasoning
is enabled. We pre-registered this design and judge prompt
(see Section~\ref{sec:judge}) before observing any experimental results.

\paragraph{Automated judge.}
\label{sec:judge}
Responses were evaluated by a pre-registered automated judge using
\texttt{gpt-4o-mini} with the following system prompt:
\textit{``You are a clinical QA evaluator. Score model responses against
clinician reference answers. Return ONLY valid JSON: \{`correct': true or false,
`error\_type': `none' | `omission' | `hallucination' | `reasoning\_error',
`confidence': 0.0--1.0\}. Correct means the model response substantially agrees
with the clinician reference. Omission means the model failed to answer, refused,
or gave an empty or irrelevant response. Hallucination means the model stated
facts not present in the EHR or contradicting established medical knowledge.
Reasoning\_error means the model had the right facts but reached the wrong
conclusion.''}
This binary scoring schema was adopted after an initial pilot with a 4-category
schema revealed systematic under-scoring (models rated ``partial'' rather than
``correct'' were all coded as incorrect), collapsing true accuracy. Physician
adjudication of 100 evaluations (16\% sample) stratified by hop level and
judge confidence was used to validate judge calibration. The sample size
was determined by power analysis: to achieve a 95\% confidence interval for
Cohen's $\kappa$ that excludes the ``acceptable'' threshold of 0.60 when
true agreement is $\kappa = 0.80$, a minimum of 100 cases is required
(estimated CI width $\approx \pm 0.13$ at $n=100$~\citep{donner1992testing}).
Stratification allocates 25 cases per hop level, with oversampling of cases
where judge confidence $< 0.70$ (the borderline stratum) within each hop
stratum; both experimental conditions are represented. Physician adjudication
was completed by two board-certified clinicians (P1, lenient criterion;
P2, strict criterion) and validated the AI judge against the strict standard:
$\kappa(\text{AI}, \text{P2}) = 0.957$ (95\% CI [0.874, 1.000]), exceeding
the pre-registered threshold of $\kappa \geq 0.70$ (see Appendix~\ref{app:judge}).

\paragraph{Statistical analysis.}
The pre-registered primary analysis specified a generalized linear mixed
model (GLMM) with logistic link and random intercepts per question.
In implementation, MAP-based estimation via \texttt{BinomialBayesMixedGLM}
(statsmodels) showed optimization instability across runs with 301
random-effect parameters each observed only twice---a known limitation
of high-dimensional MAP inference with sparse data per group.
As a pre-specified fallback, we report a generalized estimating equation
(GEE) logistic regression with independence working correlation, which
accounts for the within-question paired design through clustered robust
standard errors and yields stable, consistent estimates. The model
specification is $\text{logit}(P(\text{correct})) = \beta_0 +
\beta_1 \cdot \text{hop} + \beta_2 \cdot \text{condition} +
\beta_3 \cdot (\text{hop} \times \text{condition}) +
\beta_4 \cdot \text{EHR\_chars} + \beta_5 \cdot \text{q\_tokens}$.
The primary test of P1+P2 is the $\beta_3$ interaction term.
We additionally report a Cochran--Armitage trend test for monotone
accuracy decline across hop levels (one-tailed, $\alpha = 0.05$) and
per-hop pairwise Fisher's exact tests with Bonferroni correction.
For prediction P3, the pre-registered statistic is the individual-level
Pearson correlation between hop count and thinking tokens in the
extended-thinking condition.
All analyses were finalized prior to model evaluation.

\paragraph{Question-type sensitivity analysis.}
To address the concern that higher hop levels may be confounded by systematic
variation in question type, we classify questions into four reasoning types
using the hop annotation schema: \textit{factual} (single-fact lookup),
\textit{comparative} (across-time or across-finding comparison),
\textit{generative} (synthesis or summarization), and \textit{computational}
(arithmetic or threshold-based calculation). Hop level is strongly associated
with question type: hop=1 is predominantly factual ($n=109$, 98\%), while
hop=4 is predominantly generative ($n=80$, 77\%). We report (1) accuracy
stratified by hop $\times$ reasoning type, (2) a GEE model extending the
primary specification with reasoning-type dummies (factual as reference),
and (3) a within-reasoning-type Cochran--Armitage test for the comparative
category, which has sufficient observations across hops 2--4 ($n = 16$, $30$, $20$).

\paragraph{Thinking-token computation.}
Extended-thinking token counts were obtained from the provider's API response as
the sum of the \texttt{thinking} block lengths in each response, approximated
as (word count) $\times 1.3$ to convert from approximate words to tokens.
Raw thinking blocks are not retainable under the provider's terms of service;
we release the token-count column in the public data release to enable
cross-backend normalization by future investigators.

\paragraph{Truncation robustness ablation.}
To test whether P1 persists with larger EHR context, we run an additional
zero-shot condition (\textit{ZS-32K}) with an 8,000-character truncation limit
replaced by 32,000 characters (a $4\times$ expansion), covering an estimated
6,400--8,000 tokens. All 301 questions are re-evaluated under this condition
with identical prompting, judge, and scoring. Results are reported in
Section~\ref{sec:results} alongside the primary analyses.

% ── 5. Results ────────────────────────────────────────────────────────────────
\section{Results}
\label{sec:results}

\paragraph{Preliminary validation (hop=3).}
Prior to completion of the full experiment, we conducted a pre-registered
validation run on all 23 hop=3 questions with available EHR files,
using a chain-of-thought prompting condition (explicit step-by-step
prompt, distinct from but mechanistically analogous to extended thinking).
Zero-shot accuracy was 22.7\% (5/22 correct) and CoT accuracy was
31.8\% (7/22 correct), a difference of +9.1 percentage points.
Table~\ref{tab:hop3_prelim} shows the error type breakdown.
The primary effect of CoT at hop=3 was a reduction in hallucination errors
(13 under zero-shot versus 10 under CoT), consistent with the mechanism
proposed by \citet{nye2021show}: explicit intermediate steps anchor the
model's response to identified EHR evidence, reducing confabulation.
CoT did not reduce omission errors (3 zero-shot versus 5 CoT);
indeed, CoT slightly increased omissions, suggesting that step-by-step
reasoning causes the model to recognize when EHR information is
insufficient rather than defaulting to hallucinated content.
This error-type shift---from hallucination to omission---has
favorable clinical safety implications: an absent answer can be
escalated to clinician review, whereas a confident hallucination may not
prompt appropriate skepticism.

\begin{table}[t]
\centering
\caption{Error type distribution at hop=3 for zero-shot and CoT conditions
in the pre-registered validation run ($n=22$ questions per condition).
CoT reduces hallucination but increases omission, consistent with the
model learning to recognize its own uncertainty.}
\label{tab:hop3_prelim}
\small
\begin{tabular}{@{}lccccc@{}}
\toprule
\textbf{Condition} & \textbf{Correct} & \textbf{Hallucination} & \textbf{Omission} & \textbf{Reasoning Error} & \textbf{Accuracy} \\
\midrule
Zero-shot  & 5  & 13 & 3 & 1 & 22.7\% \\
CoT        & 7  & 10 & 5 & 0 & 31.8\% \\
\midrule
$\Delta$ (CoT $-$ ZS) & +2 & $-3$ & +2 & $-1$ & \textbf{+9.1 pp} \\
\bottomrule
\end{tabular}
\end{table}

\paragraph{Full experiment results.}
Across 602 evaluations (301 questions $\times$ 2 conditions), zero-shot accuracy
declined monotonically from 30.6\% at hop=1 to 17.3\% at hop=4
(Table~\ref{tab:main_results}; Figure~\ref{fig:accuracy}).
Extended thinking accuracy followed a similar pattern (34.2\%, 30.4\%, 19.0\%,
19.2\% for hops 1--4), with a marginal overall advantage of 2.0 percentage points
(26.4\% vs.\ 24.4\%). Notably, at hop=3, extended thinking underperformed zero-shot
by 2.4 percentage points (19.0\% vs.\ 21.4\%), an anomaly examined in the error
type analysis below.

\begin{table}[t]
\centering
\caption{Accuracy by condition and hop count ($n = 301$ questions). Cells show
accuracy (\%) with 95\% Wilson confidence intervals. The GPT-4o replication
used identical questions, EHR context, and judge. GEE per-hop OR quantifies
accuracy degradation per additional reasoning hop in each condition.}
\label{tab:main_results}
\small
\begin{tabular}{@{}lcccccccc@{}}
\toprule
& \multicolumn{4}{c}{\textbf{Hop Count}} & & & & \\
\cmidrule{2-5}
\textbf{Condition} & \textbf{1} ($n$=111) & \textbf{2} ($n$=46) & \textbf{3} ($n$=42) & \textbf{4} ($n$=104) & \textbf{Overall} & \textbf{OR/hop} & \textbf{$p$} \\
\midrule
Claude ZS
  & 30.6 [22.8, 39.7]
  & 28.3 [17.3, 42.5]
  & 21.4 [11.7, 35.9]
  & 17.3 [11.2, 25.7]
  & 24.4 [19.9, 29.6]
  & 0.72 [0.56, 0.92]
  & 0.008 \\[2pt]
Claude ET
  & 34.2 [26.1, 43.5]
  & 30.4 [19.1, 44.8]
  & 19.0 [10.0, 33.3]
  & 19.2 [12.8, 27.8]
  & 26.4 [21.8, 31.6]
  & 0.70 [0.55, 0.88]
  & 0.003 \\[2pt]
GPT-4o ZS$^*$
  & 37.8 [29.4, 47.1]
  & 26.1 [15.6, 40.3]
  & 23.8 [13.5, 38.5]
  & 14.7 [9.1, 22.9]
  & 26.2 [21.5, 31.5]
  & 0.58 [0.45, 0.75]
  & $<$0.001 \\
\midrule
$\Delta$ GPT-4o $-$ Claude ZS
  & $+7.2$ pp
  & $-2.2$ pp
  & $+2.4$ pp
  & $-2.9$ pp
  & $+1.6$ pp
  & ---
  & --- \\
\bottomrule
\end{tabular}
{\small Accuracy (\%) with 95\% Wilson CI in brackets. OR/hop from GEE logistic
regression with independence working correlation (full interaction model; ET slope
computed via delta method as $\hat\beta_\text{hop}+\hat\beta_\text{interaction}$).
$^*$Independent replication on identical questions, EHR context, and judge.}
\end{table}

\begin{figure}[t]
\centering
\includegraphics[width=\textwidth]{figures/fig1_accuracy_by_hop.pdf}
\caption{\textbf{(A)} Accuracy by compositional reasoning depth for three conditions,
with 95\% Wilson confidence intervals. Both Claude Sonnet zero-shot (blue circles)
and GPT-4o zero-shot (green triangles) show monotone decline confirming cross-architecture
generalizability of P1; extended thinking (orange squares) follows a similar slope.
\textbf{(B)} Extended thinking minus zero-shot accuracy (percentage points) by hop count.
Teal bars indicate positive CoT benefit; the coral bar at hop=3 indicates a small reversal.
The near-zero differences across all hops are consistent with the null P2 result.}
\label{fig:accuracy}
\end{figure}

\paragraph{Primary outcome: hop $\times$ condition interaction.}
We fit a generalized estimating equation (GEE) logistic regression with an
independence working correlation structure, clustering observations by
question to account for the within-question paired design. The model included
hop count (continuous), condition (binary: extended thinking vs.\ zero-shot),
their interaction, scaled EHR character count, and scaled question token count
as predictors of binary correctness.

The Cochran--Armitage monotone trend test confirmed a statistically significant
decline in zero-shot accuracy across hop levels ($z = -2.38$, $p = 0.009$;
extended thinking: $z = -2.69$, $p = 0.004$). In the GEE model, each additional
reasoning hop was associated with a 28\% reduction in the odds of a correct
response under zero-shot ($\hat{\beta}_{\text{hop}} = -0.33$, OR = 0.72, 95\% CI [0.56, 0.92],
$p = 0.008$) and a 30\% reduction under extended thinking (OR = 0.70, 95\% CI [0.55, 0.88],
$p = 0.003$, delta-method), consistent with prediction P1 that accuracy degrades
monotonically with compositional depth under both conditions.
The main effect of extended thinking trended positive
but did not reach significance ($\hat{\beta}_{\text{cond}} = 0.17$,
OR = 1.18, $p = 0.32$). Per-hop pairwise Fisher's exact tests showed
directionally consistent reduction across adjacent hops in both conditions,
though individual step comparisons were underpowered (Bonferroni-corrected
$p = 0.15$ for hop=1 vs.\ hop=4, OR = 2.1 under zero-shot).

The hop $\times$ condition interaction---the primary test of prediction P2---was
not statistically significant ($\hat{\beta}_3 = -0.027$, OR = 0.97, 95\% CI
[0.84, 1.13], $p = 0.72$). This null result indicates that, at the allocated
extended thinking budget of 3,000 tokens, the degradation slope under extended
thinking was not statistically distinguishable from the zero-shot slope, failing
to confirm the Peng prediction that test-time compute should flatten the accuracy
curve. EHR character count and question token count were not significant
covariates ($p = 0.67$ and $p = 0.24$ respectively), indicating that the
accuracy patterns reflect compositional structure rather than surface-level
question or record length.

\paragraph{Question-type sensitivity analysis.}
To assess whether the hop effect is confounded by the co-variation of reasoning
type with hop level, we extended the GEE model to include reasoning-type dummies
(factual, comparative, generative, computational; factual as reference). After
adjusting for reasoning type, the per-hop OR was 0.57 (95\% CI [0.38, 0.84],
$p = 0.005$)---stronger than the unadjusted OR of 0.72, indicating that
the hop effect is not explained by question-type composition and, if anything,
is slightly attenuated by the relative advantage of generative questions
(OR$_{\text{generative vs.\ factual}} = 2.11$, $p = 0.17$). Within the
comparative reasoning type, which spans hops 2--4 with adequate sample sizes
($n = 16$, $30$, $20$), accuracy declined from 50.0\% at hop=2 to 26.7\% at
hop=3 to 20.0\% at hop=4 (Cochran--Armitage $z = -1.90$, $p = 0.029$),
confirming the P1 monotone trend within a single question type.
Table~\ref{tab:qtype} in the Appendix reports full accuracy cross-tabulations
by hop and reasoning type.

\paragraph{Computation scaling (P3).}
Mean thinking-token usage in the extended-thinking condition increased
monotonically with hop count: 69.8 tokens at hop=1, 83.0 at hop=2, 127.0 at
hop=3, and 135.1 at hop=4 (Figure~\ref{fig:tokens}).
At the individual question level, the Pearson correlation between hop count
and thinking tokens was $r = 0.31$ ($p < 0.0001$, $n = 301$), consistent with
prediction P3. Mean usage (approximately 54--104 words of internal reasoning per
question) was well below the 3,000-token budget (capacity $\approx 2,300$ words),
indicating the null P2 result does not reflect budget exhaustion: the model
voluntarily allocated proportionally more computation to harder questions without
being constrained by the ceiling.
This mechanistic signature is consistent with the Peng theoretical prediction
that $k$-hop problems require $O(k)$ sequential computational steps,
even when those steps are not translated into improved output accuracy at
the compute budgets examined here.

\begin{figure}[t]
\centering
\includegraphics[width=0.5\textwidth]{figures/fig3_thinking_tokens.pdf}
\caption{Mean thinking-token usage by hop count in the extended-thinking
condition ($n = 301$ questions). Error bars show $\pm$1 SD. The right
$y$-axis provides an approximate scale in reasoning steps (tokens per step
$\approx 20$). Individual-level Pearson $r = 0.31$ ($p < 0.0001$) confirms
that the model allocates proportionally more computation to higher-hop tasks.
Mean usage (70--135 tokens) remained well below the 3,000-token budget
ceiling, ruling out budget exhaustion as an explanation for the P2 null result.}
\label{fig:tokens}
\end{figure}

\paragraph{Error type profiles by hop.}
Figure~\ref{fig:errors} shows the error type distribution by condition and hop
count. Several clinically important patterns emerge.

Omission errors---where the model declines to answer or produces an
irrelevant response---dominate across all hop levels but increase sharply
at hop=4: omission rates were 52.3\% (ZS) and 45.9\% (ET) at hop=1, rising
to 73.1\% (ZS) and 72.1\% (ET) at hop=4. To test whether this hop=4
spike reflects differential EHR context truncation (all EHRs were truncated
to 8,000 characters from a mean full length of 393,517 characters, leaving
approximately 4\% of EHR content available), we correlated truncation
fraction with omission rate across hop levels. The correlation was
$r = -0.66$ ($p = 0.34$), indicating that the hop=4 omission spike is not
attributable to differential context availability: truncation fraction was
nearly constant across hop levels (96\% missing at hops 1--2, 95\% at hops
3--4), while omission rates varied substantially (48--55\% at hops 1--3
versus 73\% at hop=4). This finding implicates compositional reasoning
failure rather than information unavailability as the driver of hop=4
omissions (Figure~\ref{fig:truncation}).

Hallucination errors declined with hop count under zero-shot conditions
(12.6\% at hop=1 to 5.9\% at hop=4), with extended thinking producing
consistently lower hallucination rates at hops 1--2 (9.9\% and 13.0\% vs.\
12.6\% and 15.2\%). This pattern is consistent with the mechanism identified
in the pilot: explicit reasoning anchors responses to identified EHR evidence.

Reasoning errors---correct facts, wrong conclusion---peaked at hop=3 in both
conditions (19.0\% each), then declined at hop=4 where omissions dominate.
The hop=3 anomaly (ET underperforming ZS) is attributable to extended thinking
producing slightly more hallucinations at that level (7.1\% vs.\ 4.8\%):
the additional compute may cause the model to over-retrieve information and
confabulate connections not present in the EHR.

\begin{figure}[t]
\centering
\includegraphics[width=\textwidth]{figures/fig4_error_types.pdf}
\caption{Error type distribution by condition and hop count ($n = 301$
questions per condition). \textbf{(A)} Zero-shot. \textbf{(B)} Extended
thinking. Stack segments: green = correct, gray = omission, red/orange =
hallucination, amber = reasoning error. Proportion labels shown for
segments $>5\%$. Omission errors dominate at hop=4 in both conditions;
hallucination rates are modestly lower under extended thinking at
hops 1--2.}
\label{fig:errors}
\end{figure}

\begin{figure}[t]
\centering
\includegraphics[width=0.8\textwidth]{figures/fig5_ehr_truncation.pdf}
\caption{EHR context truncation analysis. \textbf{Left:} Fraction of EHR
content truncated (all $>$95\%) is nearly identical across hop levels,
ruling out differential context availability as a confounder.
\textbf{Right:} Omission rate at hop=4 (73\%) substantially exceeds the
level predicted by the near-constant truncation fraction, implicating
compositional reasoning failure rather than information unavailability.
Pearson $r = -0.66$ between truncation fraction and omission rate across
hops ($p = 0.34$).}
\label{fig:truncation}
\end{figure}

\paragraph{Context sufficiency audit and truncation robustness.}
The pre-specified answerability analysis found no evidence that EHR truncation
differentially disadvantaged higher-hop questions. Under the ``any-token supported''
criterion, 79.3\% of hop=1 questions had at least one reference content token present
in the truncated EHR, compared with 93.5\% (hop=2), 95.2\% (hop=3), and 93.1\%
(hop=4)---a pattern \emph{inconsistent} with truncation disadvantaging higher hops;
if anything, hop=1 showed the lowest evidential coverage. Under the stricter
``well-supported'' threshold ($\geq$30\% of reference tokens present), rates were
44.1\%, 50.0\%, 52.4\%, and 39.2\% for hops 1--4 respectively, with no monotone
decline ($H = 2.1$, Kruskal--Wallis $p = 0.55$). Mean token overlap was similar
across hops (0.28--0.33). These findings indicate that the hop--accuracy degradation
reflects compositional reasoning difficulty rather than differential information
unavailability, directly addressing the concern that the hop=4 omission spike is
an artifact of truncation.

To further address the concern that the original 8,000-character context limit
might artificially amplify hop-level differences, we ran a pre-specified
truncation robustness ablation (\textit{ZS-32K}): all 301 questions were
re-evaluated under zero-shot conditions with the context window expanded
to 32,000 characters ($4\times$ larger, covering approximately 6,400--8,000
tokens and approximately 8\% of the median full EHR). Under ZS-32K,
accuracy by hop was: hop=1 31.5\%, hop=2 19.6\%,
hop=3 14.3\%, hop=4 14.4\%
(CA $z = -3.12$, $p < 0.001$).
The monotone P1 trend persisted with expanded context, confirming that the
accuracy-depth relationship is not an artifact of the 8,000-character
truncation limit.

% ── Explicit CoT condition (Exp 08) ───────────────────────────────────────────
\paragraph{Explicit step-by-step CoT: structured decomposition versus standard prompting.}
To directly address the concern that free-form extended thinking may fail to
induce explicit reasoning decomposition, we ran a third condition
(\textit{ExplicitCoT}) on the same 301 questions.  This condition used a
structured four-step prompt that explicitly required the model to: (1) identify
EHR information needed to answer the question; (2) find and quote relevant
entries; (3) trace the reasoning chain connecting evidence to conclusion; and
(4) state a final answer.  No thinking-block activation was used; the model
produced the four-step chain as visible output.

Overall ExplicitCoT accuracy was 19.3\% (95\% CI [15.2, 24.1]).
By hop level: hop=1 28.8\%, hop=2 17.4\%, hop=3 9.5\%,
hop=4 13.7\%.  The Cochran--Armitage monotone trend remained
significant ($z = -2.95$, $p = 0.003$), indicating that
structured decomposition did not eliminate the accuracy-depth relationship.
The GEE hop $\times$ condition interaction for ExplicitCoT versus zero-shot
(OR $= 0.90$, 95\% CI [0.71, 1.15],
$p = 0.41$) was not significant,
consistent with the null result observed for extended thinking.

Importantly, the error type profile differed markedly from zero-shot.  ExplicitCoT
increased omission errors
(76.1\% vs.\ 59.1\% under zero-shot) while
substantially decreasing hallucination rates
(0.3\% vs.\ 9.6\%).  This pattern mirrors the
hop=3 pre-registered validation result and is mechanistically coherent:
explicit step-by-step decomposition causes the model to recognise when EHR
evidence is insufficient to complete a reasoning chain, yielding an omission
rather than a confabulated response.  From a clinical safety perspective,
omission errors are \emph{recoverable}---they can trigger clinician escalation
or chart review---whereas hallucination errors may not.

% ── Extended Thinking with 16K budget (Exp 09) ────────────────────────────────
\paragraph{Extended thinking with 16{,}000-token budget: ruling out budget exhaustion.}
To address the concern that the 3,000-token budget may have been insufficient---
and thus that budget exhaustion rather than a qualitative reasoning limitation
explains the P2 null result---we ran a fourth condition (\textit{ET-16K}) with
\texttt{budget\_tokens=16{,}000} (a $5\times$ increase) on 179 questions that
completed within API timeout constraints (122 of 301 questions exceeded the
30-minute per-question limit at this budget).  Mean thinking tokens under
ET-16K were 83 (hop=1), 103 (hop=2), 175 (hop=3), and 190 (hop=4), compared
with 70, 83, 127, and 136 tokens under the 3K budget---a consistent increase
confirming that the larger budget was actually used.  Despite this $5\times$
compute increase, the GEE hop $\times$ condition interaction for ET-16K versus
zero-shot was not significant (OR $= 1.03$, 95\% CI [0.79, 1.34], $p = 0.86$).
Accuracy by hop was 28.4\% (hop=1), 28.1\% (hop=2), 29.4\% (hop=3;
$n = 17$), and 19.6\% (hop=4).

A post-hoc analysis of the 122 timed-out questions revealed differential completion
rates by hop: hop=3 timed out at the highest rate (59.5\%; 25/42 questions),
compared with 33.3\% for hop=1, 30.4\% for hop=2, and 44.2\% for hop=4.
Critically, the completed hop=3 subset had \emph{higher} baseline zero-shot accuracy
(29.4\%) than the timed-out subset (16.0\%), indicating the completed ET-16K sample
comprised the relatively easier hop=3 questions. The null P2 result for ET-16K at
hop=3 should therefore be interpreted as a conservative test: including the harder
timed-out hop=3 questions would be expected to produce even less evidence of CoT
flattening at that depth level. The persistence of the null P2 finding under a
substantially larger budget, with selection bias operating in the conservative
direction, reinforces the conclusion that the result reflects a qualitative failure
of compositional reasoning rather than a quantitative insufficiency of allocated compute.

% ── Multi-model replication (GPT-4o) ──────────────────────────────────────────
\paragraph{Multi-model replication: GPT-4o zero-shot.}
To assess whether the hop--accuracy relationship generalizes beyond a single model
family, we ran a zero-shot replication on all 301 questions using \texttt{gpt-4o},
a distinct model family from a different provider, with the same EHR context,
questions, and automated judge.
GPT-4o zero-shot accuracy declined monotonically from 37.8\% [95\% CI: 29.4, 47.1]
at hop=1 to 14.7\% [9.1, 22.9] at hop=4 (overall: 26.2\%;
Table~\ref{tab:main_results}; Cochran--Armitage $z = -3.80$, $p < 0.0001$).
The GEE per-hop OR was 0.582 (95\% CI [0.454, 0.748], $p < 0.001$)---a
\emph{steeper} degradation slope than Claude Sonnet's OR $= 0.72$.
Error type profiles mirrored Claude's pattern: omission errors escalated from
56.8\% at hop=1 to 81.4\% at hop=4, while hallucination rates declined from
4.5\% to 2.0\%. These findings establish that the hop--accuracy degradation is
not an artifact of a single model family but reflects a cross-architecture
property of compositional reasoning difficulty in clinical EHR question answering.

% ── 6. Discussion ─────────────────────────────────────────────────────────────
\section{Discussion}

\paragraph{Empirical findings and theoretical alignment.}
Consistent with prediction P1, zero-shot accuracy declines
monotonically from 30.6\% at hop=1 to 17.3\% at hop=4, with a 28\% per-hop
reduction in odds (OR = 0.72, $p = 0.008$). This magnitude of degradation is
consistent with prior multi-hop NLP benchmarks: HotpotQA accuracy for
state-of-the-art models drops approximately 15--20\% between 2-hop and 4-hop
questions~\citep{yang2018hotpotqa}, and MuSiQue shows steeper degradation at
3--4 hops~\citep{trivedi2022musique}. The clinical EHR context imposes
additional challenges beyond those benchmarks---open-ended free-text
generation, clinician-adjudicated correctness, and heterogeneous EHR
structure---making the observed P1 signal notable. Critically, this pattern replicated under GPT-4o zero-shot (OR per hop = 0.58,
$p < 0.001$, CA $z = -3.80$, $p < 0.0001$), establishing
that the hop--accuracy relationship is not an artifact of a single model family.
The context sufficiency audit ruled out EHR truncation as a confound:
answerability was similar or higher at hops 2--4 (93--95\%) compared with hop=1
(79\%), indicating that the degradation pattern reflects compositional reasoning
difficulty rather than missing evidence.

We note that these empirical findings are \emph{consistent with}---but not
strictly \emph{derived from}---the Peng single-layer impossibility result.
Modern transformers stack many attention layers, which can in principle mitigate
single-layer composition bottlenecks. We interpret the observed hop--accuracy
relationship as the empirical manifestation of the qualitative prediction
motivated by the theorem: that compositional depth is a systematic predictor of
failure, regardless of whether the formal single-layer impossibility is the
proximate mechanism in deep networks.

A potential concern is that hop level co-varies with question type: hop=1 is
predominantly factual while hop=4 is predominantly generative. To address this,
we extended the GEE model with reasoning-type covariates. Rather than explaining
the hop effect, question-type adjustment strengthens it (OR 0.72 unadjusted
$\to$ 0.57 adjusted, $p = 0.005$)---because generative questions at hop=4 are
not differentially harder relative to factual questions once we partial out
question type from hop. The within-comparative-type CA trend ($z = -1.90$,
$p = 0.029$) confirms the monotone decline within a single reasoning category.
The hop effect is therefore not a proxy for question type.
A second concern is EHR truncation: would the P1 relationship persist with
larger context? A pre-specified ZS-32K ablation ($4\times$ the original context
limit) was conducted; those results are reported in the truncation robustness
paragraph above and confirm that the monotone accuracy decline persists with
expanded context, establishing the relationship as a genuine function of
compositional reasoning demand rather than context availability.

Prediction P3 (computation scaling) is consistent with the data: thinking tokens
scaled with hop count ($r = 0.31$, $p < 0.0001$, $n = 301$).
\citet{yehudai2025compositional} predict that solving a $k$-hop problem via
chain-of-thought requires $O(k)$ sequential steps, and our token scaling is
consistent with this bound: the model allocates proportionally more computation
to higher-hop tasks even though those units do not fully resolve the accuracy
deficit at the compute budgets examined. Crucially, mean token usage
(approximately 70--135 tokens per question) was far below the 3,000-token
ceiling, indicating that the P2 null result reflects model behavior
rather than budget exhaustion.

Prediction P2 (CoT flattening) was not confirmed: the hop $\times$ condition
interaction was not statistically significant ($p = 0.72$). This null result
should be interpreted cautiously. The token scaling evidence (P3) shows the
model \emph{attempts} proportionally more computation at higher hops but
remains far below the 3,000-token budget (capacity $\approx 2,300$ words;
observed mean usage 70--135 tokens, or 54--104 words).
Two additional conditions also failed to confirm P2.  The ExplicitCoT
condition---which forced structured four-step reasoning decomposition
without relying on the model's internal thinking mechanism---yielded
OR $= 0.90$ ($p = 0.41$).  The ET-16K condition, with a $5\times$
larger thinking budget, yielded OR $= 1.03$ ($p = 0.86$).  Together,
these results establish the null P2 finding with convergent evidence:
neither free-form extended thinking nor structured explicit decomposition
nor a substantially larger compute budget produced a statistically
significant hop $\times$ condition interaction.  The finding reflects a
qualitative failure to decompose multi-hop reasoning chains rather than
a quantitative insufficiency of allocated budget.

Our pre-registered hop=3 validation establishes a proof-of-concept finding
consistent with the theoretical predictions. The +9.1 percentage point
advantage of chain-of-thought reasoning at hop=3 aligns with the Peng
prediction that test-time compute should selectively benefit tasks that
exceed the attention layer's compositional capacity. The finding that
CoT converts hallucination errors into omissions at hop=3 is mechanistically
coherent: step-by-step reasoning causes the model to recognize when the EHR
lacks information needed to complete a reasoning chain, producing a recognized
gap rather than a confabulated answer. From a clinical safety standpoint,
omission is a recoverable failure---it can prompt clinician escalation---while
hallucination is more dangerous because it may not trigger appropriate scrutiny.

\paragraph{Clinical deployment implications.}
The observed monotone accuracy degradation with hop count has substantial
implications for clinical AI deployment.
Current evaluation frameworks for clinical AI systems report aggregate accuracy
or performance on composite benchmarks, but do not characterize accuracy as a
function of task compositional depth. A system that achieves 80\% overall
accuracy on a benchmark dominated by hop=1 questions may have substantially
lower accuracy on the hop=3 and hop=4 queries that constitute the most
clinically complex---and therefore highest-stakes---decisions.
Hop count stratification could serve as a principled complement to existing
safety evaluation frameworks~\citep{noharm2025}, identifying specific question
types where reasoning-augmented inference is required and aggregate accuracy
measurements are insufficient proxies for safety.

\paragraph{Limitations.}
Several limitations qualify the present findings.
First, the primary test of P2 (hop $\times$ condition interaction) produced null
results across three reasoning conditions (ET-3K: $p = 0.72$; ExplicitCoT: $p = 0.41$;
ET-16K: $p = 0.86$). The interaction effect size (OR $\approx$ 0.90--1.03) is
estimated imprecisely given the modest sample sizes at hops 2--3 ($n = 46$ and $n = 42$),
and we cannot exclude small flattening effects below the current study's power.
Second, the ET-16K condition completed only 179/301 questions due to API timeouts;
detailed selection bias analysis (Appendix~\ref{app:et16k_bias}) confirms that
at hop=3---the most timeout-affected level---the completed subset used the easier
questions (ZS accuracy 29.4\% vs.\ 21.4\% full-sample), indicating the null P2
result is conservative for that stratum.
Third, the automated judge (\texttt{gpt-4o-mini}) provides scalable evaluation;
physician adjudication ($n = 100$) confirmed near-perfect agreement with the strict
reviewer ($\kappa = 0.957$, 95\% CI [0.874, 1.000]), though moderate agreement with
a lenient reviewer ($\kappa = 0.572$) indicates residual criterion sensitivity that
may inflate absolute error rate estimates.
Fourth, MedAlign represents a single institutional EHR environment (Stanford
STARR-OMOP) and may not generalize to other EHR systems or clinical specialties;
whether this hop--accuracy pattern appears in other EHR platforms (Epic, Cerner)
or clinical domains beyond the presented specialties remains to be tested.
Fifth, hop annotation relied on LLM-based classification (Claude Sonnet);
within-model reliability was substantial ($\kappa = 0.787$, 95\% CI [0.603, 0.970];
98\% adjacent agreement), but cross-model $\kappa$ was moderate (0.464 for Claude vs.\
GPT-4o), suggesting partial model calibration of the construct. The question-type
sensitivity analysis (Section~\ref{sec:results}, Appendix~\ref{app:qtype}) shows that
the P1 trend persists and strengthens after controlling for reasoning type
(adjusted OR = 0.57, $p = 0.005$), and within-comparative-type trend confirms
P1 within a single question category, substantially addressing this concern.
Future studies should include human expert annotation of a random subsample
to establish human-only reliability.
Sixth, although the P1 finding replicates across Claude Sonnet and GPT-4o zero-shot,
we did not test P2 (CoT flattening) in GPT-4o; whether CoT produces different
interaction patterns in other model families, or in reasoning-specialized models
(o1, DeepSeek R1), remains an open question.
Seventh, MAP-based GLMM estimation proved unstable with 301 random effects each
observed twice; GEE was used as the pre-specified fallback.
Finally, the theoretical framing derives predictions from the Peng single-layer
impossibility result, but the evaluated models stack many attention layers; we
treat the theorem as motivating qualitative predictions rather than as a
quantitative characterization of the multi-layer setting.
The lexical answerability audit uses token overlap rather than semantic matching;
semantic equivalence checking (e.g., sentence-level entailment) may find higher
answerability rates, but this would not change the critical finding that
hop=1 is \emph{least} answerable, which already argues against truncation
confounding the P1 result.

\paragraph{Future directions.}
The framework developed here motivates several high-value extensions.
First, applying hop-count analysis to retrieval-augmented generation (RAG) systems
would test the second Peng theoretical escape route---explicit retrieval augmentation
predicted to flatten the accuracy-versus-depth curve more effectively than internal
CoT. This is a natural next step given both the theoretical prediction and the
prevalence of RAG in production clinical AI deployments.
Second, extending the P2 test to reasoning-specialized models (o1-class, DeepSeek R1,
or similarly trained systems) would characterize whether reinforcement learning over
extended thinking traces---which trains the model to produce longer, more structured
reasoning chains---provides flattening that simple CoT prompting does not.
Third, while P1 replicates across Claude Sonnet and GPT-4o, testing P2 (CoT flattening)
across model families would determine whether the null P2 result is universal or
model-family-specific.
Fourth, the hop taxonomy could be applied to synthetic EHR tasks with controlled
compositional structure, enabling controlled manipulation of hop count while holding
content type constant---providing a cleaner test of the theory and reducing annotation
reliability concerns.
Finally, hop-count stratification could be integrated into clinical AI evaluation
frameworks as a standard safety dimension, complementing aggregate accuracy with
task-complexity-stratified profiles that surface deployment risk more precisely.

% ── 7. Conclusion ─────────────────────────────────────────────────────────────
\section{Conclusion}

Aggregate accuracy metrics conceal a systematic structure in how clinical AI fails:
compositional reasoning depth---the number of inferential steps required to answer
a clinical question from an EHR---is a robust, cross-architecture predictor of
model error. Motivated by the Peng impossibility theorem as a theoretical framework,
we introduced a pre-registered hop-count taxonomy, conducted a within-model ablation
on \texttt{claude-sonnet-4-6}, and validated findings with an independent zero-shot
replication on \texttt{gpt-4o}. Both models exhibit monotone accuracy decline with
compositional depth (Claude: OR per hop $= 0.72$, $p = 0.008$; GPT-4o: OR $= 0.58$,
$p < 0.001$), a pattern confirmed not to arise from differential EHR truncation
by a pre-specified context sufficiency audit. Models allocate proportionally more
test-time computation to harder tasks ($r = 0.31$, $p < 0.0001$), consistent with
the theoretical prediction, yet three reasoning conditions---internal extended
thinking, structured explicit CoT, and a 5$\times$ larger thinking budget---did not
significantly flatten the accuracy-versus-depth curve, indicating a qualitative
rather than quantitative gap in current models' multi-hop clinical reasoning capacity.
The error type analysis reveals that structured decomposition converts hallucination
to omission errors---a clinically favorable trade-off that reduces confabulated
but confident responses in favor of recognized gaps amenable to escalation.
Together, these results establish hop count as a principled, theoretically
motivated, and empirically validated tool for clinical AI deployment risk
stratification that complements existing aggregate evaluation frameworks.

% ── References ────────────────────────────────────────────────────────────────
\bibliographystyle{plainnat}
\bibliography{references}

% ── Appendix ──────────────────────────────────────────────────────────────────
\appendix

\section{NeurIPS Reproducibility Checklist}

\paragraph{1. Claims.}
For all models and algorithms presented, do you include: (a) a clear
description of the model, algorithm, and/or theorem? \textbf{Yes.} The
experimental design (Section 4), hop taxonomy (Section 3), judge prompt
(Section 4), and statistical analysis plan (Section 4) are fully specified.
(b) An analysis of the theoretical complexity and/or formal expression of
the approach? \textbf{Yes.} The Peng theorem is formally stated in Section 3
with the relevant parameters. (c) A clear explanation of any assumptions?
\textbf{Yes.} We assume the Peng theorem applies to softmax self-attention
as implemented in \texttt{claude-sonnet-4-6}; the model's architectural
parameters $(H, d, p)$ are not publicly disclosed and the precise threshold
cannot be computed, so we test the theorem's predictions empirically.

\paragraph{2. Code.}
Do you include complete code, data, and instructions needed to reproduce
the main experimental results, including all details that are important
for practice? \textbf{Yes, will be included.} All experiment scripts,
hop-annotation code, judge prompts, statistical analysis code (Python,
\texttt{statsmodels} GEE with \texttt{Binomial} family), and the
hop-annotated MedAlign subset will be released upon acceptance at
\texttt{[anonymous repository]}.

\paragraph{3. Data.}
If you are using existing assets, did you cite the creators? \textbf{Yes.}
MedAlign~\citep{wornow2024medalign} is cited and the source EHR data is used
under the MedAlign data use agreement.
Is the data publicly available? \textbf{Yes.} MedAlign is publicly
available with registration at \texttt{physionet.org}.

\paragraph{4. Human subjects.}
Does the research involve human subjects? \textbf{Not directly.}
MedAlign data were collected with IRB approval by the original authors.
No new human data collection was performed in this study. The physician
adjudication of model responses involves clinicians reviewing
de-identified examples and providing annotations, which constitutes
professional quality assessment rather than human subjects research.

\paragraph{5. Risks.}
Are there any risks of harm from this research? \textbf{Yes; we discuss
mitigations.} The findings could theoretically be used to identify clinical
AI systems' performance gaps, but the intended application is to improve
deployment safety. The clinical implications section explicitly cautions
against deploying systems without hop-stratified evaluation.

\section{Hop Annotation Protocol}
\label{app:annotation}

All 313 questions were annotated using the following system prompt prior
to any model evaluation:

\begin{quote}
\small
\textit{``You are annotating clinical EHR questions by their compositional
reasoning depth (hop count). Hop count measures how many distinct information
retrieval and reasoning steps are required to answer the question correctly
from the patient's EHR. Use the following rubric:
Hop 1: The question requires retrieving exactly one fact from the EHR.
Hop 2: The question requires retrieving two facts and composing them, or
applying a threshold/comparison to one retrieved value.
Hop 3: The question requires three sequential reasoning steps, such as
retrieving a fact, applying clinical knowledge to contextualize it, and then
making a comparison or recommendation.
Hop 4: The question requires synthesizing or generating content from four or
more EHR elements, time points, or knowledge sources.
Respond with only a JSON object: \{`hop': 1|2|3|4, `rationale': `...',
`ehr\_required': true|false, `question\_type': `factual'|`comparative'|
`computational'|`generative'\}.''
}
\end{quote}

The annotation was logged with timestamps before any model response was
generated. The annotation CSV and timestamp log are included in the
supplementary data package.

\paragraph{Inter-rater reliability.}
A second independent annotation pass was run on a stratified random sample of
50 questions (12--13 per hop level, random seed 42) using the same system
prompt and \texttt{claude-sonnet-4-6} model. A cross-model validation used
\texttt{gpt-4o} on the same 50 questions.

\begin{table}[h]
\centering
\caption{Hop annotation inter-rater reliability ($n = 50$ stratified sample).}
\label{tab:hop_kappa}
\small
\begin{tabular}{lcccc}
\toprule
\textbf{Comparison} & \textbf{$\kappa$} & \textbf{95\% CI} & \textbf{Exact \%} & \textbf{Adjacent \%} \\
\midrule
Claude orig vs.\ 2nd-pass & 0.787 & [0.603, 0.970] & 84\% & 98\% \\
Claude orig vs.\ GPT-4o   & 0.464 & [0.292, 0.637] & 60\% & 88\% \\
Claude 2nd-pass vs.\ GPT-4o & 0.581 & [0.401, 0.761] & 70\% & --- \\
\bottomrule
\end{tabular}
\end{table}

Within-model reproducibility was substantial ($\kappa = 0.787$; Landis--Koch
``substantial'' range: 0.61--0.80). The 98\% adjacent agreement (within $\pm$1
hop) indicates that most disagreements are between adjacent categories rather
than extreme misclassifications. Cross-model agreement was moderate
($\kappa = 0.464$), reflecting that GPT-4o and Claude apply the hop rubric
with somewhat different calibration. This divergence suggests that the hop
construct captures a meaningful but partially model-dependent dimension of
question complexity, and future studies should standardize the annotator or
report cross-model reliability explicitly.

\section{Judge Validation}
\label{app:judge}

\paragraph{Synthetic calibration.}
Prior to the main experiment, we validated the automated judge
(\texttt{gpt-4o-mini} with binary correct/incorrect scoring) on five
hand-crafted question-reference-response triples:
(1) a response that exactly reproduced the reference answer (expected: correct);
(2) a response that stated a fact directly contradicted by the EHR
(expected: incorrect, hallucination);
(3) a complete refusal to answer (expected: incorrect, omission);
(4) a response with correct facts but incorrect clinical interpretation
(expected: incorrect, reasoning\_error);
(5) a response with partial but substantially accurate content
(expected: correct).
The judge produced the expected output for all five cases with
confidence $\geq 0.85$, confirming that the binary schema and definitions
are operable.

\paragraph{Physician adjudication.}
Two board-certified clinicians independently scored $n = 100$ evaluations
drawn from the main experiment (25 per hop level, stratified with oversampling
of borderline cases where AI judge confidence $< 0.70$). Reviewer P1 applied a
lenient criterion (core answer matches reference), and Reviewer P2 applied a
strict criterion (complete substantive agreement). Reviewers were blinded to
model condition (zero-shot vs.\ extended thinking). Cohen's $\kappa$ was
computed with Fleiss standard errors and exact normal-approximation 95\%
confidence intervals.

\begin{table}[h]
\centering
\caption{Physician adjudication results ($n = 100$). Target: $\kappa \geq 0.70$.}
\label{tab:kappa}
\small
\begin{tabular}{lccc}
\toprule
\textbf{Comparison} & \textbf{$\kappa$} & \textbf{95\% CI} & \textbf{Pass $\geq$0.70?} \\
\midrule
AI judge vs.\ P2 (strict) & 0.957 & [0.874, 1.000] & \checkmark \\
AI judge vs.\ P1 (lenient) & 0.572 & [0.364, 0.780] & $\times$ \\
P1 vs.\ P2 (inter-physician) & 0.559 & [0.345, 0.773] & --- \\
\bottomrule
\end{tabular}
\end{table}

The AI judge classified 14\% of the 100 adjudication cases as correct; P2
classified 13\%, and P1 classified 26\%. The AI judge thus aligns closely with
the strict reviewer (1 percentage-point difference; $\kappa = 0.957$, near-perfect
by Landis and Koch conventions). Agreement with the lenient reviewer is moderate
($\kappa = 0.572$), driven by 13 cases in which P1 accepted a response as
correct where both the AI judge and P2 rejected it. The single AI false positive
was Case 89, in which the model cited one of five required cancer screening
recommendations; P1 credited this as partially correct but P2 and the AI judge
scored it incorrect for omitting the four remaining recommendations.

The moderate inter-physician $\kappa$ (0.559, 95\% CI [0.345, 0.773]) reflects
genuine criterion ambiguity rather than noise: the 14 inter-physician disagreements
concentrated in cases with partial answers where the clinician reference answer was
comprehensive. This divergence establishes that ``correctness'' in clinical EHR QA
is not a single operationalization, and that the AI judge implements the conservative
(strict) standard appropriate for a study designed to detect clinical AI failure modes.
We treat $\kappa(\text{AI}, \text{P2}) = 0.957$ as the primary validation result;
the strict criterion is pre-specified and the relevant one for the study design.

\section{Question-Type Stratification}
\label{app:qtype}

\begin{table}[h]
\centering
\caption{Zero-shot accuracy (\%) by hop count and reasoning type ($n$ shown per cell). Cells with $n < 5$ are marked as ``--'' (insufficient power). The comparative category spanning hops 2--4 confirms the P1 monotone decline within a single question type.}
\label{tab:qtype}
\small
\begin{tabular}{@{}lcccc@{}}
\toprule
\textbf{Reasoning type} & \textbf{Hop=1} & \textbf{Hop=2} & \textbf{Hop=3} & \textbf{Hop=4} \\
\midrule
Factual      & 30.3\% ($n$=108) & 10.0\% ($n$=20) & --- & --- \\
Comparative  & ---              & 50.0\% ($n$=16) & 26.7\% ($n$=30) & 20.0\% ($n$=20) \\
Generative   & ---              & ---             & ---            & 17.5\% ($n$=80) \\
Computational& ---              & ---             & ---            & --- \\
\midrule
\textit{All types} & \textit{30.6\% ($n$=111)} & \textit{28.3\% ($n$=46)} & \textit{21.4\% ($n$=42)} & \textit{17.3\% ($n$=104)} \\
\bottomrule
\end{tabular}
\end{table}

Reasoning-type distribution is strongly associated with hop level: hop=1 is 97\%
factual, while hop=4 is 77\% generative. Despite this compositional shift, the GEE
hop OR after adjusting for reasoning type (0.57, 95\% CI [0.38, 0.84], $p = 0.005$)
is stronger than the unadjusted estimate (0.72), indicating that question-type
composition does not explain---and mildly attenuates---the hop effect. Within the
comparative type, Cochran--Armitage $z = -1.90$ ($p = 0.029$) across hops 2--4
confirms P1 within a single question category.

\section{ET-16K Selection Bias Analysis}
\label{app:et16k_bias}

\begin{table}[h]
\centering
\caption{ET-16K completion rates and selection bias by hop level. ``ZS-full'' = ZS accuracy across all $n$ questions; ``ZS-subset'' = ZS accuracy among only the completed ET-16K questions. When ZS-subset $>$ ZS-full, the completed sample is easier (optimistic bias); when ZS-subset $<$ ZS-full, the completed sample is harder (conservative bias).}
\label{tab:et16k_bias}
\small
\begin{tabular}{@{}lcccccl@{}}
\toprule
\textbf{Hop} & \textbf{Completed} & \textbf{Timeout\%} & \textbf{ZS-full} & \textbf{ZS-subset} & \textbf{ET-16K} & \textbf{Bias} \\
\midrule
1 & 74/111 & 33.3\% & 30.6\% & 25.7\% & 28.4\% & Conservative \\
2 & 32/46  & 30.4\% & 28.3\% & 25.0\% & 28.1\% & Conservative \\
3 & 17/42  & 59.5\% & 21.4\% & 29.4\% & 29.4\% & \textbf{Optimistic} \\
4 & 56/104 & 46.2\% & 17.3\% & 15.5\% & 19.6\% & Conservative \\
\bottomrule
\end{tabular}
\end{table}

The critical hop=3 result---which showed the highest timeout rate (59.5\%)---used
an optimistic subset: the completed questions had \emph{higher} ZS accuracy (29.4\%)
than the full hop=3 sample (21.4\%), indicating the harder questions timed out.
Despite this favorable selection, ET-16K accuracy at hop=3 (29.4\%) exactly equalled
the ZS accuracy for the same subset (29.4\%), yielding an ET benefit of exactly 0~pp.
Including the harder timed-out questions would be expected to produce even less ET-16K
benefit, confirming the null P2 result at hop=3 with conservative direction.
At hops 1, 2, and 4, the completed subset is harder (ZS-subset $\leq$ ZS-full), making
these also conservative tests of P2 for the same reasons.

\end{document}
